\section*{子代数}

\begin{enumerate}
\def\labelenumi{\arabic{enumi})}
\setcounter{enumi}{12}
\item
  若 \(P \subset R\)，置 \(g^{P} = \bigoplus_{\alpha \in P} g^{\alpha}\)，\(h_{P} = \sum_{\alpha \in P} kH_{\alpha}\)。设 \(P \subset R\)，\(h'\) 为 h 的向量子空间，\(a = h' \oplus g^{P}\)。则 a 是 g 的子代数当且仅当 P 是 R 的闭子集且 \(h'\) 含有 \(h_{P \cap (-P)}\)；a 在 g 中约化当且仅当 P = -P；a 可解当且仅当 \(P \cap (-P) = \varnothing\)。
\item
  令 \(\mathbf{P}\) 为 \(\mathbf{R}\) 的闭子集，\(\mathfrak{b} = \mathfrak{h} \oplus \mathfrak{g}^{\mathrm{P}}\)。下列条件等价：
\end{enumerate}

\begin{enumerate}
\def\labelenumi{(\roman{enumi})}
\item
  \(\mathfrak{b}\) 是 \(\mathfrak{g}\) 的极大可解子代数；
\item
  \(\mathrm{P} \cap (-\mathrm{P}) = \varnothing\) 且 \(\mathrm{P} \cup (-\mathrm{P}) = \mathrm{R}\)；
\item
  存在一个房 \(\mathbf{C}\)（属于 \(\mathbf{R}\)）使 \(\mathrm{P} = \mathrm{R}_{+}(\mathrm{C})\)（参见第VI章§1第6节）。
\end{enumerate}

\((\mathfrak{g},\mathfrak{h})\) 的 Borel 子代数是指 g 的一个含 h 且满足上述条件的子代数。g 的子代数 b 若存在 g 的分裂嘉当子代数 \(h'\) 使 b 是 \((\mathfrak{g},\mathfrak{h}')\) 的 Borel 子代数，则称 b 为 g 的 Borel 子代数；若 k 代数闭，这等价于说 b 是 g 的极大可解子代数。

令 \(\mathfrak{b} = \mathfrak{h} \oplus \mathfrak{g}^{\mathrm{R}_{+}(C)}\) 为 \((\mathfrak{g}, \mathfrak{h})\) 的 Borel 子代数。\(\mathfrak{b}\) 的最大幂零理想是 \([\mathfrak{b}, \mathfrak{b}] = \mathfrak{g}^{\mathrm{R}_{+}(C)}\)。令 \(B\) 为 \(R\) 的与 \(C\) 相伴的基；代数 \([\mathfrak{b}, \mathfrak{b}]\) 由诸 \(\mathfrak{g}^{\alpha}\)（\(\alpha \in B\)）生成。

若 \(\mathfrak{b},\mathfrak{b}'\) 是 \(\mathfrak{g}\) 的 Borel 子代数，则存在 \(\mathfrak{g}\) 的嘉当子代数含于 \(\mathfrak{b} \cap \mathfrak{b}'\)；这样的子代数是分裂的。

\begin{enumerate}
\def\labelenumi{\arabic{enumi})}
\setcounter{enumi}{14}
\tightlist
\item
  令 \(\mathbf{P}\) 为 \(\mathbf{R}\) 的闭子集，\(\mathfrak{p} = \mathfrak{h} \oplus \mathfrak{g}^{\mathrm{P}}\)。下列条件等价：
\end{enumerate}

\begin{enumerate}
\def\labelenumi{(\roman{enumi})}
\item
  p 含有 \((\mathfrak{g}, \mathfrak{h})\) 的一个 Borel 子代数；
\item
  \(P \cup (-P) = R;\)
\item
  存在 R 的房 C 使 \(\mathrm{P} \supset \mathrm{R}_{+}(\mathrm{C})\)。
\end{enumerate}

\((\mathfrak{g},\mathfrak{h})\) 的抛物子代数是指 g 的一个含 h 且满足上述条件的子代数。g 的子代数 p 若存在 g 的分裂嘉当子代数 \(h'\) 使 p 是 \((\mathfrak{g},\mathfrak{h}')\) 的抛物子代数，则称 p 是抛物的。

令 \(\mathfrak{p} = \mathfrak{h} \oplus \mathfrak{g}^{\mathrm{P}}\) 为 \((\mathfrak{g}, \mathfrak{h})\) 的抛物子代数，Q 为由 \(\alpha \in \mathrm{P}\)（满足 \(-\alpha \notin \mathrm{P}\)）组成的集合，\(\mathfrak{s} = \mathfrak{h} \oplus \mathfrak{g}^{\mathrm{P} \cap (-\mathrm{P})}\)。则 \(\mathfrak{p} = \mathfrak{s} \oplus \mathfrak{g}^{\mathrm{Q}}\)，\(\mathfrak{s}\) 在 \(\mathfrak{g}\) 中约化，且 \(\mathfrak{g}^{\mathrm{Q}}\) 是 \(\mathfrak{p}\) 的最大幂零理想兼 \(\mathfrak{p}\) 的幂零根基。\(\mathfrak{p}\) 的中心为 0。

